\documentclass[12pt,a4paper]{report}

\usepackage{graphics}
\usepackage{fullpage,epsf,graphicx}%, amstext,url} 


\def\BibTeX{{\rm B\kern-.05em{\sc i\kern-.025em b}\kern-.08em
    T\kern-.1667em\lower.7ex\hbox{E}\kern-.125emX}}

\begin{document}

\thispagestyle{empty}

%
%	This is a basic LaTeX Template for the TP/MP MSc Dissertation report

\parindent=0pt          %  Switch off indent of paragraphs 
\parskip=5pt            %  Put 5pt between each paragraph  

%	This section generates a title page
%       Edit only the sections indicated to put in the project title, and submission date

\vspace*{0.1\textheight}

\begin{center}
        \huge{\bfseries Title of my Dissertation}\\ % Replace with the title of your dissertation!
\end{center}

\medskip

\begin{center}
        \Large{Katrina Ross}\\  % Author of dissertation - replace with your name!
        \medskip
        \large{August 18, 2023}  % Submission date
\end{center}

%%% If necessary, reduce the number 0.4 below so the University Crest
%%% and the words below it fit on the page.
%%% Don't let the crest, or the wording below it, flow onto the next page!

\vspace*{0.4\textheight}

\begin{center}
        \includegraphics[width=35mm]{crest.pdf}
\end{center}

\medskip

\begin{center}

%%%
%%% Change Theoretical to Mathematical if appropriate
%%%
\large{
  MSc in Theoretical Physics\\[0.8ex]
  The University of Edinburgh\\[0.8ex]
  2023}

\end{center}

\newpage


\pagenumbering{roman}

\begin{abstract}
This is where you summarise the contents of your dissertation. It should be
at least 100 words, but not more than 250 words.
\end{abstract}

\pagenumbering{roman}

\begin{center}
\textbf{Declaration}
\end{center}

I declare that this dissertation was composed entirely by myself.

Chapters 2 and 3 provide an introduction to the subject area and a
description of previous work on this topic. They do not contain
original research.

Chapter~4 describes work that was carried out entirely by me. The results of
this chapter have been obtained previously by Professor Anne T Matta of the
University of Kinlochteuchter, but the methods used here are different
in some important (or minor) ways.

Chapters 4 through 6 contain my original work. The work described in
Chapter~4 was carried out in collaboration with Professor Carole Ann O'Malley
and her PhD student Jack O'Bean. Chapter~5 presents original work done
entirely by me.

\bigskip

State whether calculations were done using Mathematica, Maple, MATLAB,
SymPy, etc, with (or without) gamma-matrix manipulation code, master
integrals, the Super-Duper software package, etc. In other words, you
should refer to any software that you used during your project. For
example, Monte Carlo simulation packages, hydrodynamics packages,
measurement code, fitting code, tensor algebra and/or calculus packages,
Feynman diagram evaluation packages, etc.

State whether any software you used was written by you from scratch,
by your supervisor (or by whoever), or if it's a standard package.

\newpage

\begin{center}
\textbf{Personal Statement}
\end{center}

\emph{You \textbf{\emph{must}} include a Personal Statement in your
  dissertation. This should describe what you did during the project,
  and when you did it. Give an account of problems you faced and how
  you attempted to overcome them. The examples below are based on
  personal statements from MSc and MPhys projects in previous years,
  with (mostly-obvious) changes to make them anonymous. }

\subsubsection*{Example~1: an analytical project}

The project began with an introduction to the spinor-helicity
formalism in four dimensions, with my main source material being
H. Elvang's “Scattering Amplitudes in Gauge Theory and Gravity” [1]. I
read the first chapter, and acquainted myself with the formalism,
and how it worked in a practical sense.

Once I felt more comfortable with it, we moved onto the
six-dimensional spinor-helicity formalism paper, where I spent some
time gaining as strong an understanding of how the formalism worked,
and proving identities.

The next stage was to learn about the generalised unitarity procedure,
with the end goal being to use it to calculate coefficients for some
one loop integral, likely involving massive particles. Learning how
this worked took some time, and proved to be some of the most
difficult material for me to understand.

It wasn't until later that we began to consider applying what I had
learned to a Kaluza-Klein reduction, which ended up being the main
focus of the project. It mixed well with the general theme of
“extra-dimensional theory” the project began with, and allowed me to
apply all that I'd learned and prepared for so far.  The vast majority
of my remaining time was spent calculating coefficients for the scalar
box contribution to the gluon-gluon to two-Kaluza-Klein-particle
amplitude, overcoming a number of problems and errors, to finally have
human-readable, and presentable results.

During the course of the project, I met with my supervisor every week,
in order to discuss my progress and the direction I would head
next. Toward the end, the frequency of our meetings increased
somewhat, as I began to finish my calculations.

I started writing this dissertation in mid-July, and I spent the first
three weeks of August working on it full-time.

Overall, I feel that the project was a success, and I found it to be
extremely enjoyable throughout.


\subsubsection*{Example~2: a computational project}

I spent the first 2 weeks of the project reading the material
surrounding my project - mainly [1] and [2]. I also began to plan out
how I would implement the algorithms in C++, in doing this I gained an
understanding of what the main goals of the first half of my project
would be and how they could be achieved. I identified which Monte
Carlo observables would be useful to measure in these simulations.

For the next 3 weeks I implemented the standard Atlantic City
algorithm and debugged my code whilst developing analysis tools in
python. I compared the results from my simulations to the results from
[3] (for the Random Osculator) and [4] for the EvenMoreRandom
Osculator. Having obtained positive results for the Random Osculator I
started reading up on Heaviside Articulation. I examined how to
integrate a Heaviside Articulator into the simulation in order to
produce the most efficient simulation - the solution I decided on was
to use a package called HeaviArt[5].

Following this I began to integrate the Heaviside Articulator into my
code and test it against the regular algorithm. In addition to this I
ran longer simulations to verify my findings without Articulation.

In mid July I finished implementing Heaviside Articulation into my
code and began looking into how to quantify any improvement in speed
gained by this algorithm. As July progressed I started looking into
how to integrate the EvenMoreRandom Osculator into my code - this was
the most complicated part of the project, as discussed in the body of
this report. Despite much effort on my part, I couldn't get the
results produced by the new algorithm to agree with the old
ones. Following further study of the literature, and long discussions
with Jack O'Bean, it turned out that the original form of Heaviside
Articulation didn't applied to the EvenMoreRandom Osculator. With the
help of Jack and my supervisor, I then developed the new version
described in this report. I also did analytical calculations of the
four-point blue function to two orders higher than had
been published previously in the literature.

For the final parts of the summer I worked mainly on perfecting the
algorithm for the Random Osculator and implementing the EvenMoreRandom
Osculators algorithm with the improved Heaviside Articulation. The
final results were encouraging, but more work is clearly needed. To
this end, I have been awarded a studentship by the British University
of Lifelong Learning to extend this work during my PhD Studies at
the non-existent Scottish Highlands Institute of Technology in
Inveroxter.

I started writing this dissertation in mid-July, and I spent the first
three weeks of August working on it full-time.


\subsubsection{Example~3: a more mathematical project}

My first two weeks of work on the project consisted of building up a
working knowledge of the algebraic structures and techniques which
would be used in the main calculations which were to be carried out,
in particular carefully reading up on the 11-dimensional case [1, 2],
the general structure of the Spencer complex of the Poincar\'e
superalgebra, the Clifford and exterior algebras over a Lorentzian
vector space in general finite dimension and the relationship between
them, the spin group, the Lorentz algebra $so(V$) and their vector and
spinor representations. As a toy calculation, I computed the first two
Spencer cohomology groups of the Poincar\'e algebra and its relevant
subalgebras in order to find its filtered (sub)deformations, finding
that these are given by a space of algebraic curvature operators. This
calculation is vastly simplified compared to the superalgebra
calculation owing to the lack of spinor structure.

In the following two weeks, I turned to the particular case of 5
dimensions, initially study- ing the quaternionic structure of the
relevant Clifford algebra, its spinor representation and symplectic
Majorana spinors. I read up on complex and quaternionic structures and
representations, proved various identities for products of higher-rank
gamma matrices found in [5] and derived the Fierz identity and various
other useful identities involving quantities defined in terms of the
Majorana spinors. I made an initial attempt to solve the relevant
cocycle conditions, initially finding that the space of solutions was
given by 3-forms (or Hodge-dually 2-forms), but from the known spinor
connection in D = 5 supergravity we knew that there was a problem with
the solution. Working with my supervisor to fix the errors in this
calculation, finding the most efficient way of setting it out and
writing up the work done so far and some of the background material
took up the following two to three weeks.  Next, I moved on to
figuring out the fundamentals of the geometric part of the project,
understanding spin structures on manifolds, the spin lift of the
Levi-Civita connection, the spinorial Lie derivative and the
definition of the superconnection. I derived the flatness conditions
for the curvature of the superconnection and showed that these
conditions are sufficient to cause the Killing superalgebra to close
and form a Lie superalgebra. This work, along with further writeup,
took another two weeks.

The final weeks of the project were spent learning about Lorentzian
and Riemannian symmetric spaces, finding all of the the maximally
supersymmetric backgrounds, on the way learning a little about the
supersymmetric solutions of 5-dimensional supergravity, and finally
finishing the dissertation.




\newpage

\begin{center}
%\vspace*{2in}
% an acknowledgements section is completely optional but if you decide
% not to include it you should still include an empty {titlepage}
% environment as this initialises things like section and page numbering.
\textbf{Acknowledgements}
\end{center}

\emph{Put your acknowledgements here. Thanking your supervisor for
his/her help is standard practice, but it's not compulsory \ldots}

I'd like to thank my supervisor Professor Carole Ann O'Malley for
making this project possible, and her PhD student Jack O'Bean for his
patience and his detailed functional explanations of how classical
symmetries can be broken by quantum effects. Thanks also to Wally Bee
and Ken Garoo of the University of Woolloomooloo for sending me
their higher-order hopping-parameter expansions.

Finally, none of this would have been possible without financial
support from Paterson's Lane Education Committee.

\bigskip

This document has its origins in the dissertation template for the MSc
in High Performance Computing, which is apparently descended from a
template developed by Professor Charles Duncan for MSc students in
Meteorology. His acknowledgement follows:

\emph{This template has been produced with help from many former
  students who have shown different ways of doing things. Please make
  suggestions for further improvements.}

Some parts of this template were lifted unashamedly from the Edinburgh
MPhys project report guide, with little or no modification. I have no
idea who wrote the first version of that\ldots

You don't have to use \LaTeX\ for your dissertation. You can use
Microsoft Word, Apple Pages, LibreOffice or similar, but it's
\emph{much} easier to typeset equations in \LaTeX, and references look
after themselves. Whatever you use, your dissertation should have the
general structure of this template, and it should look similar --
especially the front page.

\tableofcontents
\listoftables
\listoffigures

\pagenumbering{arabic}

\chapter{Introduction}
The Introduction should contain a description of your project and the
problem you are trying to solve. It should start off at a level that
should be understandable by anyone with a degree in physics, but it
can become more technical later

Where appropriate you should include references to work that has
already been done on your topic and anything else which lets you set
your work in context.

One of the things you will need to do is to ensure that you have a
suitable list of references.  To do this you should see \cite{ref:lam}
or some other suitable reference.  Note the format of the citation used
here is the style favoured in this School.  Here is another
reference \cite{ref:bloggs} for good measure.

Alternatively, you can use \BibTeX. See later for some details on this.

You will also want to make sure you have no spelling or grammatical
mistakes. To help identify speling mistukes you can use the commands
\emph{aspell}, \emph{ispell} or \emph{spell} on most Linux/Unix
computers. See the appropriate manual pages. Remember that spelling
mistakes are not the only errors which can occur. Spelling checkers
will not find errors which are, in fact, valid words such as
\emph{there} for {\em their}, nor will they find repeated repeated
words which sometimes occur if your concentration is broken when
typing. \textbf{There is no substitute for thorough proof reading!}

Your dissertation should be no longer than 15,000 words. In terms of
pages, 30 pages are ok. 50 pages are fine. But it shouldn't be
much longer than that.


\chapter{Background theory and/or theory}

\section{The easy bits}
This is just to show how to break things into sections.

Many paragraphs in this demonstration document are here to provide some
padding so that sections last for more than one page to illustrate what
happens on subsequent pages. Note that the page numbering style is usually
different on the first page of a new chapter than on subsequent pages.

Here is a padding paragraph.  Rhubarb.  More rhubarb.  Yet more
rhubarb.  Rhubarb.  More rhubarb.  Yet more rhubarb.  Rhubarb.  More
rhubarb.  Yet more rhubarb.  Rhubarb.  More rhubarb.  Yet more
rhubarb.  Rhubarb.  More rhubarb.  Yet more rhubarb.  Rhubarb.  More
rhubarb.  Yet more rhubarb.  Rhubarb.  More rhubarb.  Yet more
rhubarb.  Rhubarb.  More rhubarb.  Yet more rhubarb.  Rhubarb.  More
rhubarb.  Yet more rhubarb.  Rhubarb. More rhubarb.  Yet more rhubarb.
Rhubarb.  More rhubarb.  Yet more rhubarb.  Rhubarb.  More rhubarb.
Yet more rhubarb. Rhubarb. More rhubarb.  Yet more rhubarb.  Rhubarb.
More rhubarb. Yet more rhubarb.  Rhubarb.  More rhubarb.  Yet more
rhubarb. Rhubarb.  More rhubarb.  Yet more rhubarb.  Rhubarb.  More
rhubarb.  Yet more rhubarb.  Rhubarb. Yet more rhubarb.
Too much rhubarb. No more rhubarb!

\section{The more difficult bits}
Some bits are hard.

You might want to include an equation here:

\begin{equation}
  \delta N_{\nu} = (\delta N_{\nu})_{ex} + (\delta N_{\nu})_{au} 
  \label{equation:delsplit}
  % note that the label is optional, but it allows you to refer to this
  % equation later.
\end{equation}


Here is another padding paragraph.  Bananas.  More bananas.  Yet more
bananas.  Bananas.  More bananas.  Yet more bananas.  Bananas.  More
bananas.  Yet more bananas.  Bananas.  More bananas.  Yet more
bananas.  Bananas.  More bananas.  Yet more bananas.  Bananas.  More
bananas.  Yet more bananas.  Bananas.  More bananas.  Yet more
bananas.  Bananas.  More bananas.  Yet more bananas.  Bananas.  More
bananas.  Yet more bananas.  Bananas.  More bananas.  Yet more
bananas.  Bananas.  More bananas.  Yet more bananas.  Bananas.  More
bananas.  Yet more bananas.  Bananas.  More bananas.  Yet more
bananas.  Bananas.  More bananas.  Yet more bananas.  Bananas.  More
bananas.  Yet more bananas.  Bananas.  More bananas.  Yet more
bananas.  Bananas.  More bananas.  Yet more bananas.  Bananas.  More
bananas.  Yet more bananas.  Bananas.  More bananas.  Yet more
bananas.  Bananas.  More bananas.  Yet more bananas.  Far too many
bananas.

\subsection{Hard bits}
You might want to include another equation or three here:

\begin{equation}
  \delta N_{\nu} = (\delta N_{\nu})_{ex} + (\delta N_{\nu})_{au} 
  \label{equation:delsplit2}
  % note that the label is optional but allows you to refer to this
  % equation later.
\end{equation}

Almost the same equation again.

\begin{equation}
  \delta P_{\nu} = (\delta P_{\nu})_{ex} + (\delta Q_{\nu})_{au} 
  \label{equation:delsplit3}
  % note that the label is optional but allows you to refer to this
  % equation later.
\end{equation}

You should use a different label for each equation.

Here is a padding paragraph.  Bananas.  More bananas.  Yet more
bananas.  Bananas.  More bananas.  Yet more bananas.  Bananas.  More
bananas.  Yet more bananas.  Bananas.  More bananas.  Yet more
bananas.  Bananas.  More bananas.  Yet more bananas.  Bananas.  More
bananas.  Yet more bananas.  Bananas.  More bananas.  Yet more
bananas.  Bananas.  More bananas.  Yet more bananas.  Bananas.  More
bananas.  Yet more bananas.  Bananas.  More bananas.  Yet more
bananas.  Bananas.  More bananas.  Yet more bananas.  Bananas.  More
bananas.  Yet more bananas.  Bananas.  More bananas.  Yet more
bananas.  Bananas.  More bananas.  Yet more bananas.  Bananas.  More
bananas.  Yet more bananas.  Bananas.  More bananas.  Yet more
bananas.  Bananas.  More bananas.  Yet more bananas.  Bananas.  More
bananas.  Yet more bananas.  Bananas.  More bananas.  Yet more
bananas.  Bananas.  More bananas.  Yet more bananas. And another
equation.

\begin{equation}
  \delta Q_{\nu} = (\delta L_{\nu})_{ex} + (\delta X_{\nu})_{au} 
  \label{equation:delsplit4}
  % note that the label is optional but allows you to refer to this
  % equation later.
\end{equation}

Here is a pudding paragraph.  Rhubarb crumble.  More rhubarb crumble.
Yet more rhubarb crumble.  Rhubarb crumble.  More rhubarb crumble.
Yet more rhubarb crumble.  Rhubarb crumble.  More rhubarb crumble.
Yet more rhubarb crumble.  Rhubarb crumble.  More rhubarb crumble.
Yet more rhubarb crumble.  Rhubarb crumble.  More rhubarb crumble.
Yet more rhubarb crumble.  Rhubarb crumble.  More rhubarb crumble.
Yet more rhubarb crumble.  Rhubarb crumble.  More rhubarb crumble.
Yet more rhubarb crumble.  Apple crumble.  More apple crumble.
Yet more apple crumble.  Apple crumble.  More apple crumble.
Yet more apple crumble.  Apple crumble.  More apple crumble.
Yet more apple crumble.  Apple crumble.  More apple crumble.
Yet more rhubarb crumble.  Rhubarb crumble.  More rhubarb crumble.
Yet more rhubarb crumble.  Rhubarb crumble.  More rhubarb crumble.
Yet more rhubarb crumble.  Rhubarb crumble.  More rhubarb crumble.
Yet more rhubarb crumble.  Rhubarb crumble.  More rhubarb crumble.
Way too much rhubarb crumble.

\subsection{Even harder bits}

You might sometimes want to include equations without numbering them.
\[
  E=mc^{2}
\]
And this might be one of the places where you might want to refer to
equation (\ref{equation:delsplit}). You will usually need to use the
\LaTeX\ command twice to make cross-references like this work properly.
The cross-reference information is stored in the \emph{.aux} file so
don't delete it.

\subsubsection{Numbering}
You can keep subdividing but eventually you get to a level where
numbering stops. This text is in a subsubsection which is not numbered
by default.


\paragraph{More on numbering:}

This text is in a paragraph which is also not numbered by default and
the ``title'' of the paragraph is not on a separate line.
If you want to increase the depth to which sections are numbered you
should see the section on setting the secnumdepth counter in the manual. 


\chapter{Design and/or development (of my project)}

This section should be written in standard scientific
language. Standard techniques in your research field should not be
written out in detail. In computational projects this section should
be used to explain the algorithms used and the layout of the
computational code. A copy of the actual code may be given in the
appendices if appropriate.

This section should emphasise the philosophy of the approach used and
detail novel techniques. However please note: this section should not
be a blow-by-blow account of what you did throughout the project. It
should not contain large detailed sections about things you tried and
found to be completely wrong! However, if you find that a technique
that was expected to work failed, that is a valid result and should be
included.

Here logical structure is particularly important, and you may find
that to maintain good structure you may have to present the
explorations/calculations/computations/whatever in a different order
from the one in which you carried them out.


You might sometimes want to include multiple equations in one place
\begin{eqnarray}
  E &=& ma^{2} \\
  E &=& mb^{2} \\
  E &=& mc^{2}
\end{eqnarray}
You might want to include multiple equations in one place without
numbering them
\begin{eqnarray*}
  E &=& ma^{2} \\
  E &=& mb^{2} \\
  E &=& mc^{2}
\end{eqnarray*}
You might want to include multiple equations in one place without
numbering \emph{all} of them
\begin{eqnarray}
  E &=& ma^{2} \nonumber \\
  E &=& mb^{2} \nonumber \\
  E &=& mc^{2}
\end{eqnarray}

You might also want to include diagrams.  The example shows the use of
the special command which allows existing pdf files to be included.
You would normally keep your figures separate from the text.  These
pictures might be images or pdf output from some program.

Here, I created a figure which is centred and stretched to 30\% of the
width of the page \verb+{0.30\hsize}+ and with the height stretched by
the same amount \verb+{!}+ to preserve the aspect ratio. If you omit
the extension (ie .eps, .ps or .pdf) on the file name then \LaTeX\ will
pick up the postscript copy whereas pdflatex will automatically pick
up the PDF version.


\begin{figure}

\begin{center}
  \resizebox{0.30\hsize}{!}{\includegraphics{crest.pdf}}
\end{center}

\caption{The coloured version of the University crest. The caption should explain exactly in some detail what is displayed in the table.}
\label{fig:eucrest}

\end{figure}

You should find the file crest.pdf on this wiki.

% note that labels do not need to include a description of the object
% they are labelling but it can be helpful, eg \label{fig:figurename}.

You can use a label on a figure to refer to it later. The university
crest is in Figure~(\ref{fig:eucrest}). Note that you should not use
phrases like ``the figure above'' or ``the following figure'' since
\LaTeX\ may move the figure relative to the text if it cannot be fitted
onto the current page. The figure on the next page is an example.

\begin{figure}

\begin{center}
  \resizebox{0.75\hsize}{!}{\includegraphics{rhubarbtruss.pdf}}
\end{center}

\caption{Ach, I thought we'd thrown out all the rhubarb. Who was she
  anyway?}
\label{fig:rhubarbtruss}

\end{figure}


\chapter{Another Chapter Title}
\section{Number of Chapters}

You may vary the number of chapters. The Introduction and Background
Theory chapter are essential, although you may choose a different
title for the latter. These two introductory chapters are usually
followed by a chapter on what you did yourself, with a title such as
Design and Development, although you can choose any title you
wish. After that, you might to have another chapter, or you may go
straight to the Results and Conclusions chapter.

After the Introduction, you are free to use any chapter titles you wish.




\chapter{Results and Analysis}

This section should detail the obtained results in a clear,
easy-to-follow manner. It is important to make clear what are original
results and what are repeats of previous calculations or computations.
Remember that long tables of numbers are just as boring to read as
they are to type-in!

Use graphs to present your results wherever practicable.

Results or computations should be presented with uncertainties
(errors), both statistical and systematic where applicable.

Be selective in what you include: half a dozen \emph{e.g.}~tables that
contain wrong data you collected while you forgot to switch on the
computer are not relevant and may mask the correct results.


\section{Some results}
Here are some results.

\subsection{More results}
When showing results you are likely to use tables and graphs. You can
create tables easily in \LaTeX.

\begin{table}[h]
\begin{center}
\begin{tabular}{||l|c|l||}
\hline
\textbf{File names} & \textbf{Satellite} & \textbf{Resolution}\\
\hline
  worldr            &  Meteosat          &   5km\\
  worldg            &  Meteosat          &   5km\\
  worldb            &  Meteosat          &   5km\\
\hline
\end{tabular}
\end{center}
\caption{This is a simple table. More complicated tables can have
  headings which pass over more than one column. The caption should
  explain exactly in some detail what is displayed in the table.}
\label{simple_table}
\end{table}

If you want to produce fancier tables than shown in Table \ref{simple_table}
refer to the \LaTeX\ manual or ask Madame La Google.

One of the simplest ways to produce simple graphs is to use gnuplot
which produces \LaTeX\  output. Graph~(\ref{fig:gnu}) was produced using
gnuplot with output designated as \LaTeX\  so that a \LaTeX\  output file is
produced which you can include directly or keep separate and refer to
using the \emph{include} command.

Another approach is to draw simple figures using \emph{xfig} which allows
you to export diagrams in \LaTeX\  picture format so that the diagram can
be included directly.

Perhaps the most robust way to include graphs is to convert them to
PostScript or PDF and include them in the same was as was done in
Figure~\ref{fig:eucrest} for the University Crest. You can usually do
this with most packages, including Microsoft ones; one trick for
producing PostScript is to print to a dummy PostScript printer.

% in practice you would probably keep this in a separate file and use
% the \include{filename} command to insert it here.

\begin{figure}
% GNUPLOT: LaTeX picture
\setlength{\unitlength}{0.240900pt}
\ifx\plotpoint\undefined\newsavebox{\plotpoint}\fi
\sbox{\plotpoint}{\rule[-0.200pt]{0.400pt}{0.400pt}}%
\begin{picture}(1500,1200)(0,0)
\font\gnuplot=cmr10 at 10pt
\gnuplot
\sbox{\plotpoint}{\rule[-0.200pt]{0.400pt}{0.400pt}}%
\put(220.0,113.0){\rule[-0.200pt]{292.934pt}{0.400pt}}
\put(220.0,113.0){\rule[-0.200pt]{0.400pt}{245.477pt}}
\put(220.0,113.0){\rule[-0.200pt]{4.818pt}{0.400pt}}
\put(198,113){\makebox(0,0)[r]{$0$}}
\put(1416.0,113.0){\rule[-0.200pt]{4.818pt}{0.400pt}}
\put(220.0,317.0){\rule[-0.200pt]{4.818pt}{0.400pt}}
\put(198,317){\makebox(0,0)[r]{$0.2$}}
\put(1416.0,317.0){\rule[-0.200pt]{4.818pt}{0.400pt}}
\put(220.0,521.0){\rule[-0.200pt]{4.818pt}{0.400pt}}
\put(198,521){\makebox(0,0)[r]{$0.4$}}
\put(1416.0,521.0){\rule[-0.200pt]{4.818pt}{0.400pt}}
\put(220.0,724.0){\rule[-0.200pt]{4.818pt}{0.400pt}}
\put(198,724){\makebox(0,0)[r]{$0.6$}}
\put(1416.0,724.0){\rule[-0.200pt]{4.818pt}{0.400pt}}
\put(220.0,928.0){\rule[-0.200pt]{4.818pt}{0.400pt}}
\put(198,928){\makebox(0,0)[r]{$0.8$}}
\put(1416.0,928.0){\rule[-0.200pt]{4.818pt}{0.400pt}}
\put(220.0,1132.0){\rule[-0.200pt]{4.818pt}{0.400pt}}
\put(198,1132){\makebox(0,0)[r]{$1$}}
\put(1416.0,1132.0){\rule[-0.200pt]{4.818pt}{0.400pt}}
\put(220.0,113.0){\rule[-0.200pt]{0.400pt}{4.818pt}}
\put(220,68){\makebox(0,0){$0$}}
\put(220.0,1112.0){\rule[-0.200pt]{0.400pt}{4.818pt}}
\put(414.0,113.0){\rule[-0.200pt]{0.400pt}{4.818pt}}
\put(414,68){\makebox(0,0){$1$}}
\put(414.0,1112.0){\rule[-0.200pt]{0.400pt}{4.818pt}}
\put(607.0,113.0){\rule[-0.200pt]{0.400pt}{4.818pt}}
\put(607,68){\makebox(0,0){$2$}}
\put(607.0,1112.0){\rule[-0.200pt]{0.400pt}{4.818pt}}
\put(801.0,113.0){\rule[-0.200pt]{0.400pt}{4.818pt}}
\put(801,68){\makebox(0,0){$3$}}
\put(801.0,1112.0){\rule[-0.200pt]{0.400pt}{4.818pt}}
\put(995.0,113.0){\rule[-0.200pt]{0.400pt}{4.818pt}}
\put(995,68){\makebox(0,0){$4$}}
\put(995.0,1112.0){\rule[-0.200pt]{0.400pt}{4.818pt}}
\put(1188.0,113.0){\rule[-0.200pt]{0.400pt}{4.818pt}}
\put(1188,68){\makebox(0,0){$5$}}
\put(1188.0,1112.0){\rule[-0.200pt]{0.400pt}{4.818pt}}
\put(1382.0,113.0){\rule[-0.200pt]{0.400pt}{4.818pt}}
\put(1382,68){\makebox(0,0){$6$}}
\put(1382.0,1112.0){\rule[-0.200pt]{0.400pt}{4.818pt}}
\put(220.0,113.0){\rule[-0.200pt]{292.934pt}{0.400pt}}
\put(1436.0,113.0){\rule[-0.200pt]{0.400pt}{245.477pt}}
\put(220.0,1132.0){\rule[-0.200pt]{292.934pt}{0.400pt}}
\put(45,622){\makebox(0,0){\shortstack{This is\\the\\$y$ axis}}}
\put(828,23){\makebox(0,0){This is the $x$ axis}}
\put(828,1177){\makebox(0,0){This is a plot of $y=\sin(x)$}}
\put(220.0,113.0){\rule[-0.200pt]{0.400pt}{245.477pt}}
\sbox{\plotpoint}{\rule[-0.500pt]{1.000pt}{1.000pt}}%
\put(1306,1067){\makebox(0,0)[r]{sin(x)}}
\multiput(1328,1067)(20.756,0.000){4}{\usebox{\plotpoint}}
\put(1394,1067){\usebox{\plotpoint}}
\put(220,113){\usebox{\plotpoint}}
\multiput(220,113)(3.768,20.411){4}{\usebox{\plotpoint}}
\multiput(232,178)(4.132,20.340){3}{\usebox{\plotpoint}}
\multiput(245,242)(3.825,20.400){3}{\usebox{\plotpoint}}
\multiput(257,306)(3.884,20.389){3}{\usebox{\plotpoint}}
\multiput(269,369)(3.944,20.377){3}{\usebox{\plotpoint}}
\multiput(281,431)(4.326,20.300){3}{\usebox{\plotpoint}}
\multiput(294,492)(4.137,20.339){3}{\usebox{\plotpoint}}
\multiput(306,551)(4.276,20.310){3}{\usebox{\plotpoint}}
\multiput(318,608)(4.693,20.218){3}{\usebox{\plotpoint}}
\multiput(331,664)(4.583,20.243){2}{\usebox{\plotpoint}}
\multiput(343,717)(4.754,20.204){3}{\usebox{\plotpoint}}
\multiput(355,768)(5.034,20.136){2}{\usebox{\plotpoint}}
\multiput(367,816)(5.760,19.940){2}{\usebox{\plotpoint}}
\multiput(380,861)(5.579,19.992){3}{\usebox{\plotpoint}}
\put(398.00,923.50){\usebox{\plotpoint}}
\multiput(404,943)(7.049,19.522){2}{\usebox{\plotpoint}}
\multiput(417,979)(7.288,19.434){2}{\usebox{\plotpoint}}
\put(433.18,1021.10){\usebox{\plotpoint}}
\multiput(441,1040)(8.982,18.712){2}{\usebox{\plotpoint}}
\put(460.41,1076.97){\usebox{\plotpoint}}
\put(471.84,1094.28){\usebox{\plotpoint}}
\put(484.84,1110.41){\usebox{\plotpoint}}
\put(500.42,1124.01){\usebox{\plotpoint}}
\multiput(503,1126)(19.159,7.983){0}{\usebox{\plotpoint}}
\put(519.48,1131.37){\usebox{\plotpoint}}
\multiput(527,1132)(20.136,-5.034){0}{\usebox{\plotpoint}}
\put(539.74,1128.60){\usebox{\plotpoint}}
\put(557.04,1117.38){\usebox{\plotpoint}}
\put(570.79,1101.95){\usebox{\plotpoint}}
\put(582.44,1084.80){\usebox{\plotpoint}}
\put(593.09,1066.99){\usebox{\plotpoint}}
\multiput(601,1053)(8.430,-18.967){2}{\usebox{\plotpoint}}
\put(619.18,1010.54){\usebox{\plotpoint}}
\multiput(625,996)(7.413,-19.387){2}{\usebox{\plotpoint}}
\multiput(638,962)(6.403,-19.743){2}{\usebox{\plotpoint}}
\multiput(650,925)(5.830,-19.920){2}{\usebox{\plotpoint}}
\multiput(662,884)(5.461,-20.024){2}{\usebox{\plotpoint}}
\multiput(674,840)(5.533,-20.004){2}{\usebox{\plotpoint}}
\multiput(687,793)(4.937,-20.160){3}{\usebox{\plotpoint}}
\multiput(699,744)(4.667,-20.224){2}{\usebox{\plotpoint}}
\multiput(711,692)(4.858,-20.179){3}{\usebox{\plotpoint}}
\multiput(724,638)(4.276,-20.310){3}{\usebox{\plotpoint}}
\multiput(736,581)(4.205,-20.325){3}{\usebox{\plotpoint}}
\multiput(748,523)(4.070,-20.352){3}{\usebox{\plotpoint}}
\multiput(760,463)(4.326,-20.300){3}{\usebox{\plotpoint}}
\multiput(773,402)(3.884,-20.389){3}{\usebox{\plotpoint}}
\multiput(785,339)(3.884,-20.389){3}{\usebox{\plotpoint}}
\multiput(797,276)(4.070,-20.352){3}{\usebox{\plotpoint}}
\multiput(810,211)(3.825,-20.400){3}{\usebox{\plotpoint}}
\multiput(822,147)(3.607,-20.440){2}{\usebox{\plotpoint}}
\put(828,113){\usebox{\plotpoint}}
\end{picture}
\caption{Simple Gnuplot example. The caption should tell the reader
  what is plotted against what, and explain in some detail the various
  sets of curves of data points. It shouldn't just say ``plot of
  results for the purple function in green gauge'' without further explanation.}
\label{fig:gnu}
\end{figure}

\section{Discussion of your results}

This section should give a picture of what you have taken out of your
project and how you can put it into context.

This section should summarise the results obtained, detail conclusions
reached, suggest future work, and changes that you would make if you
repeated the project.

\chapter{Conclusions}

This is the place to put your conclusions about your work. You can
split it into different sections if appropriate. You may want to include
a section of future work which could be carried out to continue your
research.

The conclusion section should be at least one page long, preferably 2
pages, but not much longer.

\appendix
% the appendix command just changes heading styles for appendices.

\chapter{Stuff that's too detailed}

Appendices should contain all the material which is considered too
detailed to be included in the main body of the text, but which is
important enough to be included in the thesis.

Perhaps this is a good place to mention \BibTeX.

You can do references in the simple way explained in the introduction,
or you can use \BibTeX.


\section{\BibTeX}
\label{sec:bibtex}

It is convenient to use \BibTeX\ to compile your bibliography.  First
you need to create a .bib file e.g.  you may call it ref.bib Then you
can put all your references into the file with entries such as
\begin{verbatim}
@Book{ob:bornwolf,
     author = "Born, M and Wolf, E",
     title  = "Principles of Optics",
     publisher = "Cambridge University Press",
     year = 1999,
     edition = {7th},
}

@Article{jr:ashkin,
Author = {A. Ashkin and J.M. Dziedzic and J.E. Bjorkholm and S. Chu},
Title = "Observation of a single beam gradient force optical tap for 
dielectric particles",
Journal = "Optics Letters",
Volume = 11,
Pages = "288-290",
Year = 1986}

@INPROCEEDINGS{seger,
 author = {J. Seger and H.J. Brockman},
 title = {What is bet-hedging?},
 editors={P.H. Harvey and L. Partridge},
 booktitle = {Oxford Surveys in Evolutionary Biology},
 year={1987},
 page={18},
 publisher={Oxford University Press},
 place={Oxford}}
\end{verbatim}
for a book, an article in a journal or an article in a proceedings volume
respectively.

Inside your \LaTeX\ file
you should include 
\begin{verbatim}
\bibliographystyle{unsrt}                      
and
\bibliography{ref}
\end{verbatim}
The first command determines the reference style, here plain and 
unsorted. With this referencing style 
a numerical referencing system (which is now the most
common in physics literature) is used and the numbering of references
will be the order in which they appear in the document. Alternatively, 
you could use
a customised `style file' but there is no real need.  The second
command just inputs your .bib file Note that only the references cited
in the text will appear in the bibliography so you can have spare
references in your .bib file.


You use the name you have given to an entry (e.g.
for the book example above the name is ob:bornwolf)
to cite the relevant article
by using the cite command in your \LaTeX\ file e.g. 
\begin{verbatim}
\cite{ob:bornwolf}
\end{verbatim}


\section{Producing your documents using \texttt{pdflatex}}

To use pdflatex your figures need to be in pdf format.  You can convert almost any image file to pdf using \texttt{convert}.  e.g. \texttt{convert myfigure.png myfigure.pdf}.

The first time you should type:
\begin{verbatim}
  pdflatex ProjectReport
  bibtex ProjectReport
  pdflatex ProjectReport
  pdflatex ProjectReport
\end{verbatim} 
This first time you run\texttt{pdflatex} it will produce a
\texttt{ProjectReport.aux}.  The \BibTeX\ command reads in the
bibliography file and makes the files \texttt{ProjectReport.bbl} and
\texttt{ProjectReport.blg} files.  These files are read in the next
\texttt{pdflatex} command, but you'll still have ``undefined
cross-reference'' errors which are sorted out by the last
\texttt{pdflatex} command.

Subsequently, you should only need to do one (or two)
\texttt{pdflatex}s, or \texttt{pdfbibtex} followed by
\texttt{pdflatex} twice if you change any references.

\vspace{5mm} You may also use plain \texttt{latex} instead of
\texttt{pdflatex}.  This requires you to use postscript graphics
instead of pdf.




\chapter{Stuff that won't be read by anyone}

Some people include in their thesis a lot of detail, particularly lots
of tables containing raw results, figures of intermediate results, or
computer code which no-one will ever read. You should be careful that
anything like this you include should contain some element of
uniqueness which justifies its inclusion.

\begin{thebibliography}{100}

\bibitem{ref:lam} L~Lamport. \emph{1986 \LaTeX\  User's Guide and Reference Manual.} Addison Wesley, pp~242.

\bibitem{ref:bloggs} F~Bloggs. \emph{1993 \LaTeX\  Users do it in Environments.} International~Journal of Silly Findings, Volume 22, pp~23-29.

\bibitem{ref:no-one} P~Thrower. \emph{2019 Too Much Rhubarb Crumble.} Dessertation Review Letters, Volume 2, pp~1-20.

\bibitem{ref:the-early-ones} H~McDonald, S~Simpson, S~Ross and K~Green. \emph{The Ones That Got Away.} Unclear Physics, Volume 1, pp~1-68.



\end{thebibliography}


\end{document}

